\subsection{Explicit core-free local fields with a structural-cell obstruction}

\begin{claim}{B.META.19}{Proved}{Procyclic wild fixed fields fail reconstruction in every finite transcendence degree}
Keep the relative local data $D,M_0,\Omega$ of
Claim~B.META.15, and put $P=P_{M_0}$.
Choose $s\ge1$, with $s\ge2$ if $p=2$, such that
$1+p^s\in\chi_p(P)$, and choose
$\tau\in P$ with $\chi_p(\tau)=1+p^s$.
Then
\[
 C=\overline{\langle\tau\rangle}\simeq\mathbb Z_p,
 \qquad F_\tau=\Omega^C                                      \tag{META.PC1}
\]
is a henselian mixed-characteristic field of
transcendence degree $\operatorname{trdeg}_{\mathbb Q}D$,
with residue field $\overline{\mathbb F}_p$.
Neither $F_\tau$ nor any of its finite extensions
satisfies closed-point-preserving relative etale
reconstruction with the structural isomorphism retained.
For every finite $M_1/M_0$ in the fixed local algebraic closure,
\[
       \operatorname{Core}_{G_{M_1}}(C\cap G_{M_1})=1.
                                                               \tag{META.PC2}
\]
For $M_0=\mathbb Q_p$ one can take
$\chi_p(\tau)=1+p$ if $p$ is odd and
$\chi_2(\tau)=5$ if $p=2$.
Consequently this supplies negative mixed-characteristic
families in every finite transcendence degree, outside
the positive normal-core criterion of Claim~B.META.18.
\end{claim}
\begin{proof}
We first justify the choices at every prime.
The $p$-power cyclotomic extension of $\mathbb Q_p$
is totally ramified and has infinite pro-$p$ part.
A finite local extension $M_0/\mathbb Q_p$ intersects
it in a finite extension, so the cyclotomic image of
$G_{M_0}$ is open in the cyclotomic image of
$G_{\mathbb Q_p}$.
For the inertia assertion specifically,
$I_{M_0}=I_{\mathbb Q_p}\cap G_{M_0}$ is an open
subgroup of $I_{\mathbb Q_p}$; its cyclotomic image
has finite index in the image of $I_{\mathbb Q_p}$.
It therefore contains an open subgroup of $1+p\mathbb Z_p$
(of $1+4\mathbb Z_2$ when $p=2$).
The quotient $I_{M_0}/P$ is pro-prime-to-$p$.
Hence the pro-$p$ part of that image comes from $P$:
its image in $\chi_p(I_{M_0})/\chi_p(P)$ is both pro-$p$
and pro-prime-to-$p$, and is therefore trivial.
This proves the existence of $s$ and $\tau$.
For $M_0=\mathbb Q_p$ the cyclotomic description gives
the two stated explicit image values.

The continuous homomorphism $\mathbb Z_p\to C$,
$z\mapsto\tau^z$, is surjective by construction.
Its kernel is a closed subgroup of $\mathbb Z_p$.
Every nonzero such subgroup is open, which would make
$C$ finite.
But $\chi_p(\tau)=1+p^s$ has infinite order, as is
seen by its nonzero $p$-adic logarithm.
The kernel is therefore zero and $C\simeq\mathbb Z_p$.
This avoids any possible torsion qualification at $p=2$.

Since $C\subset P\subset I_{M_0}$, the field $F_\tau$
contains the strict henselization of $D$ for the chosen
valuation.  It is algebraic over $D^h$, hence henselian.
The residue field of the strict henselization is
$\overline{\mathbb F}_p$; every residue element in the
local algebraic closure is already algebraic over
$\mathbb F_p$, so the residue field remains exactly
$\overline{\mathbb F}_p$.
The extension $F_\tau/D$ is algebraic and therefore
has the stated transcendence degree.
It is not separably closed, since its absolute
Galois group is the nontrivial group $C$.

We give the retained-structural-cell obstruction
in full, including its actual point function.
This is the necessary center condition already
appearing in the proof of B.50, now applied to an
explicit local core-free field family.
Let $F$ be any characteristic-zero field with algebraic
closure $\Omega$ and a nonidentity
$c\in Z(G_F)$.
Write $X=\mathbb A^1_F$ and
$p_*:X_{\mathrm{et}}\to BG_F$ for its structural
geometric morphism.
For a continuous $G_F$-set $S$, the action of $c^{-1}$
is a $G_F$-equivariant natural automorphism of $S$;
centrality is exactly what ensures equivariance.
Applying $p^*$ gives a natural automorphism of $p^*$.
Let $\alpha_c:p_*\xrightarrow{\sim}p_*$ be its
adjoint mate.
Thus, using the right-adjoint convention of
CHW Notation~2.2, the pair
\[
                  (\operatorname{id}_{X_{\mathrm{et}}},\alpha_c)
                                                               \tag{META.PC3}
\]
is a well-typed relative geometric automorphism:
its retained structural cell has source
$p_*\operatorname{id}$ and target $p_*$
\citep[Notation~2.2]{AN_CHW2024}.
Its underlying topological map is the identity,
so it preserves closed points and is pinned.

For an actual point $a\in\Omega=X(\Omega)$,
the relative point obtained by composing with
(META.PC3) is the point $c(a)$.
Here is an explicit check of the direction of the cell.
For any etale $V\to X$, applying $c^{-1}$ to the
coordinates gives a compatible bijection
\[
                    V_{c(a)}(\Omega)\longrightarrow V_a(\Omega).
\]
These bijections identify the two geometric fibre
functors.  On an etale object pulled back from $F$
the bijection is the action of $c^{-1}$ on its
$G_F$-set of geometric points.
It is consequently the prescribed mate of $\alpha_c$
on the base, and gives the required relative
identification, including the structural cell.
The actual point function of (META.PC3) is therefore
\[
                              a\longmapsto c(a).                \tag{META.PC4}
\]

This relative morphism cannot become the etale
topos morphism of an $F$-scheme morphism
$h:\mathbb A^1_F\to\mathbb A^1_F$, even up to a
relative $2$-isomorphism.
Indeed composing such an isomorphism with every
relative $\Omega$-point would give the same actual
point by CHW Proposition~3.6: its hypotheses hold
because $F$ is perfect, $\Omega/F$ is algebraic,
and $X$ is finite type \citep[Proposition~3.6]{AN_CHW2024}.
We would have $h(a)=c(a)$ for all $a\in\Omega$.
The morphism $h$ is a polynomial in $F[T]$.
Since $c$ fixes the infinite subset $\mathbb Q\subset F$,
that polynomial agrees with $T$ at infinitely many
arguments and is therefore $T$.
But $c$ moves some element of $\Omega$, a contradiction.
Nor can localization at universal homeomorphisms
remove the contradiction: if that localized
reconstruction property held, the characteristic-zero
comparison in CHW Corollary~2.32 would realize
this relative morphism by an ordinary polynomial
on the smooth seminormal source $X$
\citep[Corollary~2.32]{AN_CHW2024}.
This proves failure in the exact relative category,
including its relative $2$-isomorphisms.

The necessary group condition is in fact slimness:
every open subgroup of $G_F$ must have trivial center.
We justify the finite-extension transfer used here.
Suppose reconstruction holds over $F$, and let
$E/F$ be finite.
An $E$-scheme is an $F$-scheme together with its
map to the finite etale object $\operatorname{Spec}E$.
On etale topoi this is the slice construction for
the corresponding finite locally constant object
of $BG_F$; a map into that slice is equivalently
a map over $BG_F$ preserving the chosen section
of the pulled-back finite object.
Given a pinned relative morphism over $E$, forget
to $F$ and reconstruct it there, using seminormal
finite-type sources as permitted by Corollary~2.32.
The resulting $F$-scheme map preserves the chosen
finite-object section: the original relative
$E$-cell supplies this equality, and the
reconstructing relative $F$-isomorphism preserves
it upon pullback.
Thus the scheme map is an $E$-scheme map.
The universal property of the slice identifies
its relative $E$-morphism with the given one;
there is no further choice of a structural cell
once that section is specified.
Faithfulness descends in the same way, or follows
by forgetting to $F$.
This proves reconstruction over every finite $E/F$.
If an open subgroup $G_E\subset G_F$ had a
nonidentity central element, (META.PC3)--(META.PC4)
over $E$ would contradict that reconstruction.

Apply this argument to $F=F_\tau$ and $c=\tau$.
Every finite extension $F'/F_\tau$ has absolute
Galois group an open subgroup of $C$, hence
$p^m\mathbb Z_p$ for some $m\ge0$.
It again has a nonidentity central element,
so the same construction proves failure over $F'$.

Finally $C\cap G_{M_1}$ is an open subgroup of
$C$ for every finite $M_1/M_0$.
Its core in $G_{M_1}$ is abelian and normal.
The group-theoretic argument in Claim~B.META.18
proves that every nontrivial closed normal subgroup
of a finite-local-field absolute Galois group is
nonabelian.  The core must consequently be trivial,
which proves (META.PC2).
The same conclusion holds for the absolute Galois
subgroup of every finite $F'/F_\tau$.

For any $n\ge0$, choose $n$ algebraically independent
elements of $\mathbb Q_p$ over $\mathbb Q$ and let
$D$ be the field they generate over $\mathbb Q$.
Its closure is $\mathbb Q_p$, so the construction
applies with the explicit cyclotomic values stated
above and with transcendence degree $n$.
\end{proof}

\paragraph{Literature boundary.}
Asayama's Theorem~1.4 and Corollary~1.5 produce
Kummer-faithful algebraic extensions of number fields
with abelian absolute Galois group.
His Remark~1.6 records the gap in the formerly
asserted general implication from Kummer-faithfulness
to slimness \citep{AN_Asayama2026}.
This distinguishes Kummer-faithfulness from the
center obstruction proved above.
No henselian property of Asayama's examples is
used or asserted here, and no Kummer-faithfulness
claim is made for the explicit fields (META.PC1).
The remaining local subgroup gate is now restricted
to slim groups whose cores in every open
local absolute Galois group are trivial.
No priority or external human verification is claimed.
